\documentclass[12pt,aspectratio=169]{beamer}

% 自訂字體的封包
\usepackage{fontspec}

% 導入圖形與表格的封包
\usepackage{graphicx}
\usepackage{booktabs}
\usepackage{tikz}
\usetikzlibrary{arrows.meta, positioning, shapes.geometric}

% 提供進階數學環境與公式排版
\usepackage{mathtools, amsmath, amsfonts}

% 設定中文字體
\usepackage{xeCJK}
\setCJKmainfont{微軟正黑體}
\setmainfont{Times New Roman}
\XeTeXlinebreaklocale "zh"
\XeTeXlinebreakskip = 0pt plus 1pt

% 自訂字體顏色的封包
\usepackage{color}

% 圖表自動編號的封包
\usepackage{caption}
% \usepackage{siunitx}

% 可指定表格排版的封包
\usepackage{array}
\usepackage{makecell}

% 序列標號
\usepackage{enumerate}

% 設定自動編號
\setbeamertemplate{caption}[numbered]

% 自訂顏色
\definecolor{MyBlue}{RGB}{33, 73, 120}
\setbeamercolor{structure}{fg=MyBlue}
\setbeamercolor{title}{fg=MyBlue}

% (動畫)字會若隱若現的顯示
\beamersetuncovermixins{\opaqueness<1->{45}}{\opaqueness<1->{95}}

% 設定中文的標籤
\renewcommand{\figurename}{圖}
\renewcommand{\tablename}{表}

% 排版形式 
\mode<presentation>{
% 主題風格
\usetheme{Madrid}
% 主題顏色
\usecolortheme{Default}
% 主題字體
\usefonttheme{professionalfonts}
% itemize 及 item 的形狀
\setbeamertemplate{itemize items}[triangle]
\setbeamertemplate{itemize subitem}[circle]
\setbeamertemplate{itemize subsubitem}[square]

% 調整外框形式的字體大小
\setbeamerfont{headline}{size=\scriptsize}
\setbeamerfont{footline}{size=\scriptsize}
% 取消右下方的跳轉工具列
\setbeamertemplate{navigation symbols}{} 
% 自定義2：清除外框下緣但僅出頁碼
\setbeamertemplate{footline}[page number] 
% 設定頁面邊界
\setbeamersize{text margin left=0.8cm, text margin right=0.8cm}
\special{papersize=\the\paperwidth,\the\paperheight}
\providecommand{\tabularnewline}{\\}
}

\title{黃丞暘｜面試簡報}
\author{國立東華大學應用數學研究所}
\institute{應徵職位：軟體研發 / 資料分析 / 演算法 相關職缺}
\date{\today}

\begin{document}

% 封面
\begin{frame}
  \titlepage
\end{frame}

% 目錄
\begin{frame}{\textbf{目錄}}
  \tableofcontents
\end{frame}

%---------------------------------
\section{自介}

\begin{frame}{\textbf{自介}}
  \begin{itemize}
    \item 學歷背景：國立東華大學應用數學研究所碩二，中山醫學大學心理學系
    \item 具 5 年以上教學經驗，以及助教、專案助理、軟硬體維護實務經驗
    \item 工作特質：具教學溝通、問題整理與專案協作能力
  \end{itemize}
\end{frame}

%---------------------------------
\section{技能}

\begin{frame}{\textbf{技能}}
  \begin{block}{硬實力}
    \begin{itemize}
      \item Python: 資料分析、統計檢定、機器學習、軟體研發（學習中）
      \item R / SQL: 基本語法、PostgreSQL ( 學習中 )
      \item AI 工具：ChatGPT 輔助程式開發與學習
      \item 英文能力：多益 815（2025）
      \item 其他：GitHub、Latex、文書處理
    \end{itemize}
  \end{block}
  \begin{block}{軟實力}
    \begin{itemize}
      \item 溝通表達
      \item 快速學習
      \item 問題整理
      \item 專案協作
    \end{itemize}
  \end{block}
\end{frame}

%---------------------------------
\section{經驗與成果}

\begin{frame}{\textbf{經驗與成果}}
  \begin{block}{教育大數據校內競賽（金牌）}
    \begin{itemize}
      \item 負責資料清理、資料分析與視覺化
      \item 協助團隊規劃進度與整合分析結果
    \end{itemize}
  \end{block}

  \begin{block}{統計諮詢課程}
    \begin{itemize}
      \item 學習需求釐清、問題定義、Gage R\&R 與 SPC 基礎概念
      \item 建立統計方法連結實務問題的初步能力
    \end{itemize}
  \end{block}

  \begin{block}{個人小型專案}
    \begin{itemize}
      \item 練習 Python 後端開發與 PostgreSQL 串接
      \item 嘗試建置口袋式記帳系統，累積資料庫與系統整合經驗
    \end{itemize}
  \end{block}
\end{frame}

%---------------------------------
\section{論文研究}

\begin{frame}{\textbf{論文研究}}
  \begin{itemize}
    \item 研究主題：半導體製程資料中的模型局部決策因子分析
    \item 核心問題：資料複雜，機器學習模型的內部決策過程不易理解
    \item 我的方法：以簡單線性迴歸刻劃模型，分析特定樣本附近的重要變數
    \item 研究價值：協助工程人員理解複雜製程資料中的關鍵因子
  \end{itemize}
\end{frame}
%---------------------------------
\section{未來規劃}

\begin{frame}{\textbf{未來規劃}}
  \begin{block}{就職前準備}
    \begin{itemize}
      \item 持續以 AI 工具輔助學習 Python/Go 後端開發與 SQL 應用
      \item 強化英文技術閱讀與基本溝通能力
    \end{itemize}
  \end{block}
  \begin{block}{未來 3 年目標}
    \begin{itemize}
      \item 累積專案經驗，具備參與海外專案與外派支援的能力
      \item 強化英文溝通能力，提升技術會議與工作交流的表達能力
      \item 希望在工作中結合技術能力與國際合作經驗，拓展職涯發展方向
    \end{itemize}
  \end{block}
\end{frame}

\begin{frame}
  \centering
  \Huge{\textbf{謝謝聆聽}}
\end{frame}

\end{document}
